\chapter{safeguard our privacCan a Faraday’s cage in today’s digital age y?}
\vskip -15pt


\begin{center}
{\large\uppercase{Kumar Ankit}}\footnote{Kumar Ankit is based at New Delhi and is gainfully employed with Vedanta Limited, a mining company as their Head of Legal, External Affairs \& Company Secretary for Iron Ore and Ports division. He completed his graduation in law from Gujarat National Law University, Gandhinagar and is a fellow member of the Institute of Company Secretaries of India besides being a member of Association of Certified Fraud Examiner, Texas, US. He has been a regular speaker on legal aspects of data and internet governance at various forums. The views expressed here are his personal views.}, 


\end{center}

\vskip 2.3cm



\vfill
\newpage

\begin{multicols}{2}

\section{INTRODUCTION}

We all grew up learning from experiments and observations, making sure that we learn and understand things with interaction In India, we are more familiar with cricket and football than lawn-tennis or golf, or we all learn our mother tongue by being exposed to it. Imagine me having exposure to only indoor variety of games, would I be develop a taste for outdoor games like cricket or football. This learning by exposure is also equally easily applicable on food, religion, language and other habits, which define a person, his thoughts and his habits. The larger question is if we can weave a new habit without having an exposure. How do we know what we don’t know when our choices to know are restricted? 

In the increasing digital market of trans-boundary content, goods and services the choices for us are often made by a super intelligent, biased algorithm, backed by immense computing power specifically designed to detect tastes of sports, food, music, living lifestyle, depending on the way we have used our smartphones, connected devices, location shares, search strings and photographs. We are aware that our browsers use cookies and trackers, have ability to smartly direct us to the algorithm chosen choices. I would like to imagine a situation where our actions, emotions and urge to know is not used against us by directing us to content which fills up the void of our search with content that has been tailored fit to make us feel happy.

This leads to a questions whether we exercise the choices we actually want to exercise or do we just go by the flow of algorithms to let it determine what we should experience? This is the question for today which has implications on how the global economy sustains and caters to its users.

\section{GENERATION OF DATA}

The modern discussion over privacy has been resultant with the increasing use of data posting by gullible users (70\% of them are below age of 35\footnote{\url{https://www2.deloitte.com/content/dam/Deloitte/in/Documents/technology-media-telecommunications/in-tmt-rise-of-on-demand-content.pdf }}) using smartphones\footnote{\url{https://webobjects.cdw.com/webobjects/media/pdf/Solutions/Networking/White-Paper-Cisco-The-Zettabyte-Era-Trends-and-Analysis.pdf }}, social media, digital records, and knowing individuals as a customer. The ‘personal data’ has been closely linked with identification of an individual either directly or indirectly. Given that all readers are using devices connected to a connected network, it is estimated that by 2020, 3.4 devices per captia will be connected to the online network collectively generating and handling $\sim$2.3 ZB\footnote{1 ZB = 1billion TBs= 1trillion GBs. Each ZB is estimated at \$100 Trillion.} of data\footnote{Supra Note}. Data is often volunteered to be disclosed like name, contact details, or is observed as during browsing or purchasing, or is stealthily observed in the background like location. It is inevitable that the future beholds that the generation of data will keep increasing in geometric progression with better access to connectivity and convenience.

It is also equally true that some of the world largest company like Alphabet generate around 85\% of their global revenues ($\sim$\$ 136 Billions) from advertisement using search platforms like google, content aggregators like YouTube, google maps and hardware. Similarly, Facebook revenue from advertisement has risen sharply. This has virtually created an oligopoly in the advertisement industry.

With such vast amount of data getting generated, it is important that ethical and acceptable practices are followed by the data collectors / controllers, or by data processors / aggregators and third party data handlers. The user (data generators) should be self-aware and explicitly informed about the way the personal, transactional or behavioral data is used, and hence a regulated practices should be followed for those who are constantly monitoring the data and purpose for which the data is retained and used. As more and more activities move online, where people will research, shop and find product recommendations, the importance of privacy and data protection will  increasingly be recognized especially when the current system for data protection is highly fragmented, with divergent regulatory approaches for global and national practices. 

Personal information is often shared by people freely, and without giving a secondary consideration. A mobile phone number is almost becoming necessary to procure anything, and the most likely reason is for the person to come back to the store. In-fact a lot of inputs like birthdate, anniversary, residence and office address, handwriting, signatures are often collected in name of the visitor log/comments books. In India, the government agencies also collect a lot of data, which at times is often used in silos and remains non-aggregated. India witnesses a high incidence of data breaches (as per a report in 2018 it was $\sim$35\% of overall data breaches of the world)\footnote{Monitor-Deloitte Report, 2018} as we currently do not have safe data security practices for both online and offline access to data. Data breach by Insiders (employees or 3^{\rm rd} party vendors) with an access to data is often a significant risk as the data collectors find in increasingly difficult to find bandwidth and resources to maintain adequate data security measures.
\begin{figure}[H]
\centering
\includegraphics[scale=1.2]{src/Figures/chap2/fig01.jpg}
\caption{}\label{fig01}
\end{figure}

\section{EXERCISE OF FREE WILL WITH CHOICES}

Choice is an interesting metaphor for will. As human beings we are all entrusted a birth right of having a free will. We are all aware that information is power, but it also now well established that information is also money.  However, the way the information is gathered and used is a matter of scrutiny. The information about choice of user is generated by web analytics and results in more online advertisement getting generated by social media optimization, or banner marketing, search engine advertisement, marketing emails, online video or local online advertisement. I would instead urge people to explore aspects of their individuality to exercise free choice.

It is equally important that people react to the online advertisements much more than the offline or yellow classified book days. Digital videos are often embedded with advertisement, which our sub-conscious minds picks up subtly as the preferable option in case an option of choice arises. Every website is beaming with native advertisement, in-feed advertisements in such ways that it is hard to differentiate between the original content and the advertised ones. The data generated by the internet users along with their preferences and behavior is the Holy Grail for most advertisers. However the reality is that Google and Facebook have more magnitude, and more ways to analyze it. 

Users of wearable devices and social media networks may not conceive of themselves as having volunteered data but their activities of use and engagement result in the generation of vast amounts of data about individual lifestyles, choices and preferences. We need to our choices are future ready in view of technological advancements and social developments and do we need Privacy Enhancement Technologies (PETs) which allow online users to protect the privacy of their Personally Identifiable Information (PII) provided to and handled by Data Governors (DGs). Further, the balance between privacy and security measures is turning to be tilting towards the security measures, but at the same time should not have anything to do with the consumer behavior.

Privacy and data protection require that information about individuals should not be automatically made available to other individuals and organizations. The dangers to privacy in an age of information can originate not only from the State but from non-State actors as well. However, what happens when you have consented to provide data to an aggregator but in turn it has been handed over to a third party processor with whom you do not have any privy of contract. In case of a data breach who is the one responsible, you, the aggregator, third party or the unknown who got access to the data due to actions or inaction by any of the above. The yardstick for determining the scope of data control should be that each person must be able to exercise a substantial degree of control over that data and its use. Data protection is legal safeguard to prevent misuse of information. The challenge is the way the information is getting generated. It is a known fact that penetration of smartphones has given unheard access of information to all. The pros of getting information are weighed against the misuse of information getting collated. Data mining processes together with knowledge discovery can be combined to create facts about individuals\footnote{R. Pildes, Why rights are not trumps: social meanings, expressive harms, and constitutionalism, 27(2) The Journal of Legal Studies (1998) at pp. 725-763}. 

My one and half year old kid who barely knows how to communicate using words and translates mumbled information by using facial expressions, laughs and cries is fond of popular online content streaming on a smart TV. He is quite happy when we play him a cartoon from a favourite channel. However, similar cartoon in a different channel does not make him happy. The issue is that he has been so accustomed to the similar make video being played in loop by the algorithm that he is not able to exercise an emotion to anything else which is available on the YouTube. As I stated earlier, he made a choice based on the limited amount of information which has been fed to him by the algorithm. If the kids can be tuned into liking or disliking things by mere repetition of similar contents based on what you have watched previously, it will lead a market domination of someone who can aptly pay for maintaining that dominion based on better probabilistic abilities. The exercise of choice is getting undermined heavily by not knowing what we don’t know.

The genX, millennial and the genZ are living with smartphones as their ‘constant companion’ and are eager to know about updates regarding their friends, celebrities, sports, actions, movies, and the things which are happening in and around the world.

We get so tempted to read all the useless information on the internet so often that we have lost the purpose for which the internet was created, and have become the commodity instead of using internet as a commodity.

\section{GENESIS AND ORIGIN OF PRIVACY}

In 1891, the American lawyers Samuel Warren and Louis Brandeis described the right to privacy in a famous article: it is the right to be let alone\footnote{Holvast J. (2009) History of Privacy. In: Matyáš V., Fischer-Hübner S., Cvrček D., Švenda P. (eds) The Future of Identity in the Information Society. Privacy and Identity 2008. IFIP Advances in Information and Communication Technology, vol 298. Springer, Berlin, Heidelberg available at\url{ https://doi.org/10.1007/978-3-642-03315-5\_2}}. Technology did not rule its subject then as now, but the fundamental questions still remains that do we have a right to be let alone. Privacy in modern times is the most valuable substance known to man and a man should do all to protect it. A brief history of the major developments in the area of privacy have been depicted for understanding the nuances of the societal ask and sovereign reaction to such asks.

\begin{itemize}

\item[{\bf 1)}] \textbf{1950 - Universal Declaration of Human Rights}

When United Nations proclaimed the draft of the Universal Declaration of Human Rights (‘UDHR’) in a general assembly on December 10, 1948 at Paris, little was known that the Article 12 of the UDHR\footnote{Amended in 1950.} would lay foundation stone for a regime on which the entire current generation would depend for their fundamental right to make and exercise choices. Article 12\footnote{Article 12 of UDHR states “No one shall be subjected to arbitrary interference with his privacy, family, home or correspondence, nor to attacks upon his honour and reputation. Everyone has the right to the protection of the law against such interference or attacks.” \url{https://www.un.org/en/universal-declaration-human-rights/ }} casted obligation / duty on the member states to guarantee freedom of privacy to its subject.

\item[{\bf 2)}] \textbf{1980 - OECD Privacy Guidelines}

Thereafter in 1980, the OECD Privacy Guidelines of 1980\footnote{Accessed from \url{https://www.oecd.org/internet/ieconomy/oecdguidelinesontheprotectionofprivacyandtransborderflowsofpersonaldata.htm on 18.11.19 at 14:33 hours.}} (Guidelines) for the first time established an agreed set of privacy principles. These guidelines were drafted under the chairmanship of The Hon. Mr. Justice M.D. Kirby, Chairman of the Australian Law Reform Commission, and for the first time they recognized the fact that participating countries should have a common interest in protecting privacy and individual liberties; and processing and trans-border flows of personal data requires the development of compatible rules and practices which do not hinder flows or create unjustified obstacles ; and that is it inevitable to avoid the free flow trans-border transfer of data as it contributes significantly to the economic and social development of the countries. These Guidelines laid principles which prescribed that the collection

\item[{\bf 1)}]

\item[{\bf 1)}]

\item[{\bf 1)}]


\end{itemize}

\end{multicols}
